\documentclass[12pt]{article}

% Idioma y codificación
\usepackage[spanish, es-tabla, es-notilde]{babel}       %es-tabla para que se titule "Tabla"
\usepackage[utf8]{inputenc}

% Márgenes
\usepackage[a4paper,top=3cm,bottom=2.5cm,left=3cm,right=3cm]{geometry}

% Comentarios de bloque
\usepackage{verbatim}

% Paquetes de links
\usepackage[hidelinks]{hyperref}    % Permite enlaces
\usepackage{url}                    % redirecciona a la web

% Más opciones para enumeraciones
\usepackage{enumitem}

% Personalizar la portada
\usepackage{titling}

% Paquetes de tablas
\usepackage{multirow}

% Para añadir el símbolo de euro
\usepackage{eurosym}

% Paquete para subrayar con diferentes estilos
\usepackage[normalem]{ulem}


%------------------------------------------------------------------------

%Paquetes de figuras
\usepackage{caption}
\usepackage{subcaption} % Figuras al lado de otras
\usepackage{float}      % Poner figuras en el sitio indicado H.


% Paquetes de imágenes
\usepackage{graphicx}       % Paquete para añadir imágenes
\usepackage{transparent}    % Para manejar la opacidad de las figuras

% Paquete para usar colores
\usepackage[dvipsnames, table, xcdraw]{xcolor}
\usepackage{pagecolor}      % Para cambiar el color de la página

% Habilita tamaños de fuente mayores
\usepackage{fix-cm}

% Para los gráficos
\usepackage{tikz}
\usepackage{forest}

% Para poder situar los nodos en los grafos
\usetikzlibrary{positioning}

% Para evitar problemas entre babel y tikz
\usetikzlibrary{babel}


%------------------------------------------------------------------------

% Paquetes de matemáticas
\usepackage{mathtools, amsfonts, amssymb, mathrsfs}
\usepackage[makeroom]{cancel}     % Simplificar tachando
\usepackage{polynom}    % Divisiones y Ruffini
\usepackage{units} % Para poner fracciones diagonales con \nicefrac

\usepackage{pgfplots}   %Representar funciones
\pgfplotsset{compat=1.18}  % Versión 1.18

\usepackage{tikz-cd}    % Para usar diagramas de composiciones
\usetikzlibrary{calc}   % Para usar cálculo de coordenadas en tikz

%Definición de teoremas, etc.
\usepackage{amsthm}
%\swapnumbers   % Intercambia la posición del texto y de la numeración

\theoremstyle{plain}

\makeatletter
\@ifclassloaded{article}{
  % --- Configuración para ARTICLE ---
  \newtheorem{teo}{Teorema}[section]
  \theoremstyle{definition}
  \newtheorem{definicion}{Definición}[section]
  \newtheorem{ejercicio}{Ejercicio}[section]

  \renewcommand{\theteo}{%
    \ifnum\value{section}=0 \arabic{teo}\else \thesection.\arabic{teo}\fi
  }
  \renewcommand{\thedefinicion}{%
    \ifnum\value{section}=0 \arabic{definicion}\else \thesection.\arabic{definicion}\fi
  }
  \renewcommand{\theejercicio}{%
    \ifnum\value{section}=0 \arabic{ejercicio}\else \thesection.\arabic{ejercicio}\fi
  }
}{
  % --- Configuración para BOOK / REPORT ---
  \newtheorem{teo}{Teorema}[chapter]
  \theoremstyle{definition}
  \newtheorem{definicion}{Definición}[chapter]
  \newtheorem{ejercicio}{Ejercicio}[section]

  % Si no hay sección: Capítulo.Teorema (ej. 2.1)
  % Si hay sección:    Capítulo.Sección.Teorema (ej. 2.3.1)
  \renewcommand{\theteo}{%
    \ifnum\value{section}=0
      \thechapter.\arabic{teo}%
    \else
      \thesection.\arabic{teo}% (\thesection ya contiene el número de capítulo)
    \fi
  }
  \renewcommand{\thedefinicion}{%
    \ifnum\value{section}=0
      \thechapter.\arabic{definicion}%
    \else
      \thesection.\arabic{definicion}%
    \fi
  }
  \renewcommand{\theejercicio}{%
    \ifnum\value{section}=0
      \arabic{ejercicio}%
    \else
      \thesection.\arabic{ejercicio}% Formato por defecto: ej. 2.2.1
    \fi
  }
}
\makeatother

% Teoremas dependientes o con contador compartido
\newtheorem{coro}{Corolario}[teo]
\newtheorem{prop}[teo]{Proposición}
\newtheorem{lema}[teo]{Lema}
\newtheorem*{lema*}{Lema}

\theoremstyle{remark}
\newtheorem*{observacion}{Observación}

\theoremstyle{definition}
\newtheorem*{notacion}{Notación}
\newtheorem*{ejemplo}{Ejemplo}
\newtheorem*{ejercicio*}{Ejercicio}


% \renewcommand\qedsymbol{$\blacksquare$}         % Cambiar símbolo QED
%------------------------------------------------------------------------

% Paquetes para encabezados
\usepackage{fancyhdr}
\pagestyle{fancy}
\fancyhf{}

\newcommand{\helv}{ % Modificación tamaño de letra
\fontfamily{}\fontsize{12}{12}\selectfont}
\setlength{\headheight}{15pt} % Amplía el tamaño del índice


%\usepackage{lastpage}   % Referenciar última pag   \pageref{LastPage}
%\fancyfoot[C]{%
%  \begin{minipage}{\textwidth}
%    \centering
%    ~\\
%    \thepage\\
%    \href{https://losdeldgiim.github.io/}{\texttt{\footnotesize losdeldgiim.github.io}}
%  \end{minipage}
%}
\fancyfoot[C]{\thepage}
\fancyfoot[R]{\href{https://losdeldgiim.github.io/}{\texttt{\footnotesize losdeldgiim.github.io}}}

%------------------------------------------------------------------------

% Conseguir que no ponga "Capítulo 1". Sino solo "1."
\makeatletter
\@ifclassloaded{book}{
  \renewcommand{\chaptermark}[1]{\markboth{\thechapter.\ #1}{}} % En el encabezado
    
  \renewcommand{\@makechapterhead}[1]{%
  \vspace*{50\p@}%
  {\parindent \z@ \raggedright \normalfont
    \ifnum \c@secnumdepth >\m@ne
      \huge\bfseries \thechapter.\hspace{1em}\ignorespaces
    \fi
    \interlinepenalty\@M
    \Huge \bfseries #1\par\nobreak
    \vskip 40\p@
  }}
}
\makeatother

%------------------------------------------------------------------------
% Paquetes de cógido
\usepackage{minted}
\renewcommand\listingscaption{Código fuente}

\usepackage{fancyvrb}
% Personaliza el tamaño de los números de línea
\renewcommand{\theFancyVerbLine}{\small\arabic{FancyVerbLine}}

% Estilo para C++
\newminted{cpp}{
    frame=lines,
    framesep=2mm,
    baselinestretch=1.2,
    linenos,
    escapeinside=||
}

% para minted
\definecolor{LightGray}{rgb}{0.95,0.95,0.92}
\setminted{
    linenos=true,
    stepnumber=5,
    numberfirstline=true,
    autogobble,
    breaklines=true,
    breakautoindent=true,
    breaksymbolleft=,
    breaksymbolright=,
    breaksymbolindentleft=0pt,
    breaksymbolindentright=0pt,
    breaksymbolsepleft=0pt,
    breaksymbolsepright=0pt,
    fontsize=\footnotesize,
    bgcolor=LightGray,
    numbersep=10pt
}


\usepackage{listings} % Para incluir código desde un archivo

\renewcommand\lstlistingname{Código Fuente}
\renewcommand\lstlistlistingname{Índice de Códigos Fuente}

% Definir colores
\definecolor{vscodepurple}{rgb}{0.5,0,0.5}
\definecolor{vscodeblue}{rgb}{0,0,0.8}
\definecolor{vscodegreen}{rgb}{0,0.5,0}
\definecolor{vscodegray}{rgb}{0.5,0.5,0.5}
\definecolor{vscodebackground}{rgb}{0.97,0.97,0.97}
\definecolor{vscodelightgray}{rgb}{0.9,0.9,0.9}

% Configuración para el estilo de C similar a VSCode
\lstdefinestyle{vscode_C}{
  backgroundcolor=\color{vscodebackground},
  commentstyle=\color{vscodegreen},
  keywordstyle=\color{vscodeblue},
  numberstyle=\tiny\color{vscodegray},
  stringstyle=\color{vscodepurple},
  basicstyle=\scriptsize\ttfamily,
  breakatwhitespace=false,
  breaklines=true,
  captionpos=b,
  keepspaces=true,
  numbers=left,
  numbersep=5pt,
  showspaces=false,
  showstringspaces=false,
  showtabs=false,
  tabsize=2,
  frame=tb,
  framerule=0pt,
  aboveskip=10pt,
  belowskip=10pt,
  xleftmargin=10pt,
  xrightmargin=10pt,
  framexleftmargin=10pt,
  framexrightmargin=10pt,
  framesep=0pt,
  rulecolor=\color{vscodelightgray},
  backgroundcolor=\color{vscodebackground},
}

%------------------------------------------------------------------------
% La etiqueda IA sirve para indicar a partir de donde le pasamos el preambulo a la IA para el escaneo
% IA

% Comandos definidos
\newcommand{\bb}[1]{\mathbb{#1}}
\newcommand{\cc}[1]{\mathcal{#1}}

% I prefer the slanted \leq
\let\oldleq\leq % save them in case they're every wanted
\let\oldgeq\geq
\renewcommand{\leq}{\leqslant}
\renewcommand{\geq}{\geqslant}

% Si y solo si
\newcommand{\sii}{\iff}

% MCD y MCM
\DeclareMathOperator{\mcd}{mcd}
\DeclareMathOperator{\mcm}{mcm}

% Signo
\DeclareMathOperator{\sgn}{sgn}

% Imagen
\DeclareMathOperator{\Ima}{Im}

% Aplicación identidad
\newcommand{\id}{\mathrm{Id}}

% Letras griegas
\newcommand{\eps}{\epsilon}
\newcommand{\veps}{\varepsilon}
\newcommand{\lm}{\lambda}

\newcommand{\ol}{\overline}
\newcommand{\ul}{\underline}
\newcommand{\wt}{\widetilde}
\newcommand{\wh}{\widehat}

\let\oldvec\vec
\renewcommand{\vec}{\overrightarrow}

% Derivadas parciales
\newcommand{\del}[2]{\frac{\partial #1}{\partial #2}}
\newcommand{\Del}[3]{\frac{\partial^{#1} #2}{\partial #3^{#1}}}
\newcommand{\deld}[2]{\dfrac{\partial #1}{\partial #2}}
\newcommand{\Deld}[3]{\dfrac{\partial^{#1} #2}{\partial #3^{#1}}}


\newcommand{\AstIg}{\stackrel{(\ast)}{=}}
\newcommand{\Hop}{\stackrel{L'H\hat{o}pital}{=}}

\newcommand{\red}[1]{{\color{red}#1}} % Para integrales, destacar los cambios.

% Método de integración
\newcommand{\MetInt}[2]{
    \left[\begin{array}{c}
        #1 \\ #2
    \end{array}\right]
}

% Declarar aplicaciones
% 1. Nombre aplicación
% 2. Dominio
% 3. Codominio
% 4. Variable
% 5. Imagen de la variable
\newcommand{\Func}[5]{
    \begin{equation*}
        \begin{array}{rrll}
            \displaystyle #1:& \displaystyle  #2 & \longrightarrow & \displaystyle  #3\\
               & \displaystyle  #4 & \longmapsto & \displaystyle  #5
        \end{array}
    \end{equation*}
}

%------------------------------------------------------------------------


\newcommand{\myhref}[2]{\href{#1}{\color{blue}\uline{#2}}}


\begin{document}

    % 1. Foto de fondo
    % 2. Título
    % 3. Encabezado Izquierdo
    % 4. Color de fondo
    % 5. Coord x del titulo
    % 6. Coord y del titulo
    % 7. Fecha

    
    % 1. Foto de fondo
% 2. Título
% 3. Encabezado Izquierdo
% 4. Color de fondo
% 5. Coord x del titulo
% 6. Coord y del titulo
% 7. Fecha
% 8. Autor

\newcommand{\portada}[8]{

    \portadaBase{#1}{#2}{#3}{#4}{#5}{#6}{#7}{#8}
    \portadaBook{#1}{#2}{#3}{#4}{#5}{#6}{#7}{#8}
}

\newcommand{\portadaExamen}[8]{

    \portadaBase{#1}{#2}{#3}{#4}{#5}{#6}{#7}{#8}
    \portadaArticle{#1}{#2}{#3}{#4}{#5}{#6}{#7}{#8}
}




\newcommand{\portadaBase}[8]{

    % Tiene la portada principal y la licencia Creative Commons
    
    % 1. Foto de fondo
    % 2. Título
    % 3. Encabezado Izquierdo
    % 4. Color de fondo
    % 5. Coord x del titulo
    % 6. Coord y del titulo
    % 7. Fecha
    % 8. Autor
    
    
    \thispagestyle{empty}               % Sin encabezado ni pie de página
    \newgeometry{margin=0cm}        % Márgenes nulos para la primera página
    
    
    % Encabezado
    \fancyhead[L]{\helv #3}
    \fancyhead[R]{\helv \nouppercase{\leftmark}}
    
    
    \pagecolor{#4}        % Color de fondo para la portada
    
    \begin{figure}[p]
        \centering
        \transparent{0.3}           % Opacidad del 30% para la imagen
        
        \includegraphics[width=\paperwidth, keepaspectratio]{../../_assets/#1}
    
        \begin{tikzpicture}[remember picture, overlay]
            \node[anchor=north west, text=white, opacity=1, font=\fontsize{60}{90}\selectfont\bfseries\sffamily, align=left] at (#5, #6) {#2};
            
            \node[anchor=south east, text=white, opacity=1, font=\fontsize{12}{18}\selectfont\sffamily, align=right] at (9.7, 3) {\href{https://losdeldgiim.github.io/}{\textbf{Los Del DGIIM}, \texttt{\footnotesize losdeldgiim.github.io}}};
            
            \node[anchor=south east, text=white, opacity=1, font=\fontsize{12}{15}\selectfont\sffamily, align=right] at (9.7, 1.8) {Doble Grado en Ingeniería Informática y Matemáticas\\Universidad de Granada};
        \end{tikzpicture}
    \end{figure}
    
    
    \restoregeometry        % Restaurar márgenes normales para las páginas subsiguientes
    \pagecolor{white}       % Restaurar el color de página
    
    
    \newpage
    \thispagestyle{empty}               % Sin encabezado ni pie de página
    \begin{tikzpicture}[remember picture, overlay]
        \node[anchor=south west, inner sep=3cm] at (current page.south west) {
            \begin{minipage}{0.5\paperwidth}
                \href{https://creativecommons.org/licenses/by-nc-nd/4.0/}{
                    \includegraphics[height=2cm]{../../_assets/Licencia.png}
                }\vspace{1cm}\\
                Esta obra está bajo una
                \href{https://creativecommons.org/licenses/by-nc-nd/4.0/}{
                    Licencia Creative Commons Atribución-NoComercial-SinDerivadas 4.0 Internacional (CC BY-NC-ND 4.0).
                }\\
    
                Eres libre de compartir y redistribuir el contenido de esta obra en cualquier medio o formato, siempre y cuando des el crédito adecuado a los autores originales y no persigas fines comerciales. 
            \end{minipage}
        };
    \end{tikzpicture}
    
    
    
    % 1. Foto de fondo
    % 2. Título
    % 3. Encabezado Izquierdo
    % 4. Color de fondo
    % 5. Coord x del titulo
    % 6. Coord y del titulo
    % 7. Fecha
    % 8. Autor


}


\newcommand{\portadaBook}[8]{

    % 1. Foto de fondo
    % 2. Título
    % 3. Encabezado Izquierdo
    % 4. Color de fondo
    % 5. Coord x del titulo
    % 6. Coord y del titulo
    % 7. Fecha
    % 8. Autor

    % Personaliza el formato del título
    \pretitle{\begin{center}\bfseries\fontsize{42}{56}\selectfont}
    \posttitle{\par\end{center}\vspace{2em}}
    
    % Personaliza el formato del autor
    \preauthor{\begin{center}\Large}
    \postauthor{\par\end{center}\vfill}
    
    % Personaliza el formato de la fecha
    \predate{\begin{center}\huge}
    \postdate{\par\end{center}\vspace{2em}}
    
    \title{#2}
    \author{\href{https://losdeldgiim.github.io/}{Los Del DGIIM, \texttt{\large losdeldgiim.github.io}}
    \\ \vspace{0.5cm}#8}
    \date{Granada, #7}
    \maketitle
    
    \tableofcontents
}




\newcommand{\portadaArticle}[8]{

    % 1. Foto de fondo
    % 2. Título
    % 3. Encabezado Izquierdo
    % 4. Color de fondo
    % 5. Coord x del titulo
    % 6. Coord y del titulo
    % 7. Fecha
    % 8. Autor

    % Personaliza el formato del título
    \pretitle{\begin{center}\bfseries\fontsize{42}{56}\selectfont}
    \posttitle{\par\end{center}\vspace{2em}}
    
    % Personaliza el formato del autor
    \preauthor{\begin{center}\Large}
    \postauthor{\par\end{center}\vspace{3em}}
    
    % Personaliza el formato de la fecha
    \predate{\begin{center}\huge}
    \postdate{\par\end{center}\vspace{5em}}
    
    \title{#2}
    \author{\href{https://losdeldgiim.github.io/}{Los Del DGIIM, \texttt{\large losdeldgiim.github.io}}
    \\ \vspace{0.5cm}#8}
    \date{Granada, #7}
    \thispagestyle{empty}               % Sin encabezado ni pie de página
    \maketitle
    \vfill
}
    \portadaExamen{ffccA4.jpg}{Análisis Funcional\\Examen XVII}{Análisis Funcional. Examen XVII}{MidnightBlue}{-8}{28}{2026}{Arturo Olivares Martos}

    \begin{description}
        \item[Asignatura] Análisis Funcional.
        \item[Curso Académico] 2024/2025.
        \item[Grado] Matemáticas.
        % \item[Grupo] ---.
        \item[Profesor] Miguel Martín Suárez.
        \item[Descripción] Parcial 2.
        \item[Fecha] 4 de diciembre de 2024.
        % \item[Duración] ---.
    
    \end{description}
    \newpage


    % ------------------------------------
    
    \begin{ejercicio}
        En el espacio de Banach $X = (C\left[0, 1\right], \|\cdot\|_\infty)$ se considera el funcional lineal $\varphi$ definido por
        \[
            \varphi(f) = \int_0^1 (2t - 1) f(t) \, dt \quad \forall f \in C\left[0, 1\right].
        \]
        \begin{enumerate}[label=\alph*)]
            \item Pruebe que $\varphi$ es continuo y calcule su norma.
            \item ¿Alcanza $\varphi$ su norma? Justifique su respuesta.
            \item ¿Existe $x^{**} \in X^{**}$ con $\|x^{**}\| = 1$ tal que $x^{**}(\varphi) = \|\varphi\|$? Justifique su respuesta.
        \end{enumerate}
    \end{ejercicio}

    \begin{ejercicio}
        Rodee las afirmaciones verdaderas (que pueden ser una, dos, tres o cuatro):
        \begin{enumerate}[label=\alph*)]
            \item Sea $X$ un subespacio cerrado de $\ell_{3000}$ de dimensión infinita ¿Puede ser $X$ isomorfo a $\ell_1$?
            \begin{enumerate}[label={a\arabic*)}]
                \item Sí, porque $\ell_1 \subseteq \ell_{3000}$
                \item No, porque $\ell_1 \nsubseteq \ell_{3000}$
                \item No, porque $\ell_1$ es separable y $X$ no lo es
                \item No, porque $X^*$ es separable y $\ell_1^*$ no lo es
            \end{enumerate}
            \item Sea $\psi \colon \ell_2 \to \bb{R}$ el funcional lineal dado por $\psi(x) = \sum\limits_{n=1}^\infty \frac{x(n)}{n}$ para $x \in \ell_2$:
            \begin{enumerate}[label={b\arabic*)}]
                \item $\psi$ alcanza su norma porque $\ell_2$ es reflexivo
                \item $\psi$ alcanza su norma porque $\ell_2$ es un espacio de Hilbert
                \item $\psi$ no alcanza su norma
                \item $\psi$ no alcanza su norma porque $\ell_2$ no es completo
            \end{enumerate}
            \item Sea $Y = \bb{K}^{(\bb{N})}$ (el espacio vectorial de las sucesiones casi nulas):
            \begin{enumerate}[label={c\arabic*)}]
                \item la norma de la suma y la norma del máximo son equivalentes en $Y$
                \item la norma de la suma y la norma del máximo no son equivalentes en $Y$
                \item cualesquiera dos normas definidas en $Y$ son equivalentes
                \item no se puede definir ninguna norma en $Y$
            \end{enumerate}
            \item En un espacio normado $X$ tomamos un conjunto convexo y cerrado $C$ y un punto $x_0 \notin C$:
            \begin{enumerate}[label={d\arabic*)}]
                \item existe $f \in X^*$ tal que $\sup \Re f(C) \leq \Re f(x_0)$
                \item existe $f \in X^*$ tal que $\sup \Re f(C) < \Re f(x_0)$
                \item no se puede asegurar que $C$ se pueda separar de $x_0$ si el interior de $C$ es vacío
                \item existe un único $f \in X^*$ tal que $\sup \Re f(C) \leq \Re f(x_0)$
            \end{enumerate}
            \item Sea $Z$ un espacio normado de dimensión finita:
            \begin{enumerate}[label={e\arabic*)}]
                \item cualquier aplicación lineal que sale de $Z$ es abierta
                \item cualquier aplicación lineal que llega a $Z$ es continua
                \item cualquier aplicación lineal sobreyectiva que llega a $Z$ es abierta
                \item cualquier aplicación lineal que sale de $Z$ es continua
            \end{enumerate}
        \end{enumerate}
    \end{ejercicio}

    \newpage
    \setcounter{ejercicio}{0} % Reiniciar contador de ejercicios
    \noindent
    \textbf{Solución.}


    \begin{ejercicio}
        En el espacio de Banach $X = (C\left[0, 1\right], \|\cdot\|_\infty)$ se considera el funcional lineal $\varphi$ definido por
        \[
            \varphi(f) = \int_0^1 (2t - 1) f(t) \, dt \quad \forall f \in C\left[0, 1\right].
        \]
        \begin{enumerate}[label=\alph*)]
            \item Pruebe que $\varphi$ es continuo y calcule su norma.
            
            Veamos que $\varphi$ es continuo. Para ello:
            \begin{align*}
                |\varphi(f)| &= \left|\int_0^1 (2t - 1) f(t) \, dt\right| 
                \leq \int_0^1 |2t - 1| |f(t)| \, dt \\
                &\leq \|f\|_\infty \int_0^1 |2t - 1| \, dt
                = \|f\|_\infty \left[\int_0^{\nicefrac{1}{2}} (1 - 2t) \, dt + \int_{\nicefrac{1}{2}}^1 (2t - 1) \, dt\right] \\
                &= \|f\|_\infty \left[\left(t - t^2\right)\Big|_0^{\nicefrac{1}{2}} + \left(t^2 - t\right)\Big|_{\nicefrac{1}{2}}^1\right] = \|f\|_\infty \left[\left(\frac{1}{2} - \frac{1}{4}\right) + \left(1 - 1 - \left(\frac{1}{4} - \frac{1}{2}\right)\right)\right] \\
                &= \|f\|_\infty \left[\frac{1}{4} + \left(-\frac{1}{4} + \frac{1}{2}\right)\right] 
                = \|f\|_\infty \cdot \frac{1}{2}.
            \end{align*}

            Por lo tanto, $\varphi$ es continuo y $\|\varphi\| \leq \nicefrac{1}{2}$. Veamos ahora que $\|\varphi\| = \nicefrac{1}{2}$. La función teórica que deberíamos tomar es $f(t) = \operatorname{sign}(2t - 1)$, pero esta función no es continua. Sin embargo, podemos aproximarla por funciones continuas. Para cada $n \in \bb{N}$, definimos por comodidad:
            \begin{equation*}
                I_n = \left[\frac{1}{2} - \frac{1}{n}, \frac{1}{2} + \frac{1}{n}\right].
            \end{equation*}
            
            Definimos entonces la función continua:
            \begin{align*}
                f_n(t) = \begin{cases}
                    -1, & t \in [0, \nicefrac{1}{2} - \nicefrac{1}{n}] \\
                    1 & t \in [\nicefrac{1}{2} + \nicefrac{1}{n}, 1] \\
                    \text{lineal} & t \in I_n
                \end{cases}
            \end{align*}

            Notemos que no es necesario dar la expresión explícita de la parte lineal, ya que lo único que nos interesa es que $f_n$ sea continua y que $\|f_n\|_\infty = 1$. Estas funciones se muestran en la Figura~\ref{fig:funciones_fn}. 
            \begin{figure}[h]
                \centering
                \begin{tikzpicture}
                    \begin{axis}[
                        axis lines = middle,
                        xlabel = $t$,
                        ylabel = {$f_n(t)$},
                        ymin = -1.3, ymax = 1.3,
                        xmin = 0, xmax = 1,
                        xtick = {0, 0.5, 1},
                        ytick = {-1, 0, 1},
                        grid = both,
                        width = 11cm,
                        height = 6.5cm,
                        legend pos = south east,
                        legend cell align = left
                    ]

                    \foreach \n/\col in {3/blue, 5/red, 10/teal, 20/orange, 50/purple} {
                        \edef\temp{
                            \noexpand\addplot[color=\col, thick] coordinates {
                                (0, -1)
                                (0.5 - 1/\n, -1)
                                (0.5 + 1/\n, 1)
                                (1, 1)
                            };
                            \noexpand\addlegendentry{$n = \n$}
                        }
                        \temp
                    }

                    \end{axis}
                \end{tikzpicture}
                \caption{Gráfica de la función $f_n$ para distintos valores de $n$.}
                \label{fig:funciones_fn}
            \end{figure}

            Ahora, para cada $n \in \bb{N}$, tenemos que
            \begin{align*}
                \varphi(f_n) &= \int_0^1 (2t - 1) f_n(t) \, dt \\
                &= \int_0^1 (2t - 1)(f_n(t) - \operatorname{sign}(2t - 1)) \, dt + \int_0^1 (2t - 1) \operatorname{sign}(2t - 1) \, dt
                =\\&= \int_{I_n} (2t - 1)(f_n(t) - \operatorname{sign}(2t - 1)) \, dt + \int_0^1 |2t - 1| \, dt
            \end{align*}
            \begin{itemize}
                \item Por un lado, como hemos comprobado antes, se tiene:
                \begin{equation*}
                    \int_0^1 |2t - 1| \, dt = \frac{1}{2}.
                \end{equation*}
                \item Por otro lado, para todo $t \in I_n\subset [0, 1]$, se tiene que:
                \begin{align*}
                    |2t - 1| &\leq 1\\
                    |f_n(t) - \operatorname{sign}(2t - 1)| &\leq |f_n(t)| + |\operatorname{sign}(2t - 1)| \leq 1 + 1 = 2.
                \end{align*}

                Por lo tanto, se tiene que:
                \begin{align*}
                    \left|\int_{I_n} (2t - 1)(f_n(t) - \operatorname{sign}(2t - 1)) \, dt\right| & \leq \int_{I_n} |2t - 1| |f_n(t) - \operatorname{sign}(2t - 1)| \, dt \leq \\ & \leq 2 \int_{I_n} dt = 2 \lm(I_n) = 2 \cdot \frac{2}{n} = \frac{4}{n}.
                \end{align*}
            \end{itemize}

            Por lo tanto, para todo $n \in \bb{N}$, se tiene que:
            \begin{equation*}
                |\varphi(f_n)| \geq \frac{1}{2} - \frac{4}{n}.
            \end{equation*}

            Por tanto, se deduce que:
            \begin{equation*}
                \|\varphi\| \geq \frac{1}{2} - \frac{4}{n} \quad \forall n \in \bb{N}.
            \end{equation*}

            Tomando el límite cuando $n \to \infty$, se obtiene que:
            \begin{equation*}
                \|\varphi\| \geq \frac{1}{2}.
            \end{equation*}

            Combinando esto con la desigualdad anterior, se concluye que:
            \begin{equation*}
                \|\varphi\| = \frac{1}{2}.
            \end{equation*}



            \item ¿Alcanza $\varphi$ su norma? Justifique su respuesta.
            
            Veamos que $\varphi$ no alcanza su norma. Supongamos que existe $f_0 \in C\left[0, 1\right]$ con $\|f_0\|_\infty = 1$ tal que $|\varphi(f_0)| = \|\varphi\| = \nicefrac{1}{2}$. Entonces, se tendría que:
            \begin{align*}
                \frac{1}{2} = |\varphi(f_0)| &= \left|\int_0^1 (2t - 1) f_0(t) \,dt\right| \leq \int_0^1 |2t - 1| |f_0(t)| \, dt \leq \\ &\leq \int_0^1 |2t - 1| \|f_0\|_\infty \, dt = \int_0^1 |2t - 1| \, dt = \frac{1}{2}.
            \end{align*}

            Por lo tanto, se tendrían que cumplir todas las desigualdades anteriores como igualdades. En particular, se tendría que:
            \begin{itemize}
                \item Por un lado:
                \begin{equation*}
                    \int_0^1 |2t - 1| |f_0(t)| \, dt = \int_0^1 |2t - 1| \|f_0\|_\infty \, dt
                    \Longrightarrow
                    \int_0^1 |2t - 1| (\|f_0\|_\infty - |f_0(t)|) \, dt = 0.
                \end{equation*}

                Como se trata de una integral de una función continua y no negativa, se deduce que:
                \begin{equation*}
                    |2t - 1| (\|f_0\|_\infty - |f_0(t)|) = 0 \qquad \forall t \in [0, 1].
                \end{equation*}

                De esto se deduce que $|f_0(t)| = \|f_0\|_\infty = 1$ para todo $t \in [0, 1] \setminus \{\nicefrac{1}{2}\}$; y como $f_0$ y la norma son continuas, se deduce que:
                \begin{equation*}
                    |f_0(t)| = 1 \qquad \forall t \in [0, 1].
                \end{equation*}

                Notemos que, como $f_0$ puede tomar valores complejos, no se puede llegar a una contradicción directa.

                \item Por otro lado, se tendría que:
                \begin{equation*}
                    \left|\int_0^1 (2t - 1) f_0(t) \, dt\right| = \int_0^1 |2t - 1| |f_0(t)| \, dt
                \end{equation*}

                Como $f_0(t)$ puede ser complejo, el valor de la integral puede ser complejo. Por tanto $\exists \theta\in \bb{R}$ tal que:
                \begin{equation*}
                    \int_0^1 (2t - 1) f_0(t) \, dt = e^{i\theta} \left|\int_0^1 (2t - 1) f_0(t) \, dt\right| = e^{i\theta} \int_0^1 |2t - 1| |f_0(t)| \, dt
                \end{equation*}

                Por lo tanto, multiplicando por $e^{-i\theta}$, se deduce que:
                \begin{equation*}
                    \int_0^1 e^{-i\theta} (2t - 1) f_0(t) \, dt = \int_0^1 |2t - 1| |f_0(t)| \, dt
                \end{equation*}

                Tomando la parte real, se obtiene que:
                \begin{equation*}
                    \int_0^1 \Re\left(e^{-i\theta} (2t - 1) f_0(t)\right) \, dt = \int_0^1 |2t - 1| |f_0(t)| \, dt
                \end{equation*}

                Como $|e^{-i\theta}|=1$ y la parte real de un número complejo es menor o igual que su módulo, se deduce que:
                \begin{equation*}
                    \Re\left(e^{-i\theta} (2t - 1) f_0(t)\right) \leq |2t - 1| |f_0(t)| \qquad \forall t \in [0, 1].
                \end{equation*}

                Por lo tanto, se tendría que:
                \begin{equation*}
                    \int_0^1 \left(|2t - 1| |f_0(t)| - \Re\left(e^{-i\theta} (2t - 1) f_0(t)\right)\right) \, dt = 0
                \end{equation*}

                Como se trata de una integral de una función continua y no negativa, se deduce que:
                \begin{equation*}
                    |2t - 1| |f_0(t)| = \Re\left(e^{-i\theta} (2t - 1) f_0(t)\right) \qquad \forall t \in [0, 1].
                \end{equation*}

                De aquí, se deduce que $e^{-i\theta} (2t - 1) f_0(t)$ es un número real no negativo para todo $t \in [0, 1]$, luego:
                \begin{equation*}
                    e^{-i\theta} f_0(t) = \frac{|2t - 1| |f_0(t)|}{2t - 1}
                    \Longrightarrow
                    f_0(t) = e^{i\theta} \frac{|2t - 1|}{2t - 1} \qquad \forall t \in [0, 1] \setminus \{\nicefrac{1}{2}\}.
                \end{equation*}

                Estudiemos la continuidad en $t = \nicefrac{1}{2}$:
                \begin{align*}
                    \lim_{t \to \nicefrac{1}{2}^-} f_0(t) &= \lim_{t \to \nicefrac{1}{2}^-} e^{i\theta} \frac{|2t - 1|}{2t - 1} = \lim_{t \to \nicefrac{1}{2}^-} e^{i\theta} \frac{-(2t - 1)}{2t - 1} = -e^{i\theta},\\
                    \lim_{t \to \nicefrac{1}{2}^+} f_0(t) &= \lim_{t \to \nicefrac{1}{2}^+} e^{i\theta} \frac{|2t - 1|}{2t - 1} = \lim_{t \to \nicefrac{1}{2}^+} e^{i\theta} \frac{2t - 1}{2t - 1} = e^{i\theta}.
                \end{align*}

                Como $f_0$ es continua, se tiene que $e^{i\theta} = -e^{i\theta}$. No existe ningún número complejo con módulo 1 que cumpla esta igualdad, por lo que se llega a una contradicción. Por lo tanto, $\varphi$ no alcanza su norma.                
            \end{itemize}
            \item ¿Existe $x^{**} \in X^{**}$ con $\|x^{**}\| = 1$ tal que $x^{**}(\varphi) = \|\varphi\|$? Justifique su respuesta.
            
            Sí, existe.
            \begin{description}
                \item[Opción 1:] Por el Teorema de Hahn-Banach, como $\varphi \in X^*$ y $\varphi \neq 0$, existe $x^{**} \in X^{**}$ con $\|x^{**}\| = 1$ tal que $x^{**}(\varphi) = \|\varphi\|$.\\
                
                Aunque no se pide explícitamente, veamos por qué existe $x^{**} \in X^{**}$ con $\|x^{**}\| = 1$ tal que $x^{**}(\varphi) = \|\varphi\|$. Consideremos el subespacio $M=\bb{K}\varphi \subset X^*$ y el funcional lineal $x_0^{**} \in M^*$ definido por:
                \Func{x_0^{**}}{M}{\bb{K}}{\lambda \varphi}{\lambda \|\varphi\|}
                
                Calculamos la norma de $x_0^{**}$:
                \begin{align*}
                    \|x_0^{**}\| &= \sup_{\substack{f \in M \\ f \neq 0}} \frac{|x_0^{**}(f)|}{\|f\|} = \sup_{\substack{\lambda \in \bb{K} \\ \lambda \neq 0}} \frac{|x_0^{**}(\lambda \varphi)|}{\|\lambda \varphi\|} = \sup_{\substack{\lambda \in \bb{K} \\ \lambda \neq 0}} \frac{|\lambda| \|\varphi\|}{|\lambda| \|\varphi\|} = 1.
                \end{align*}

                Por tanto, $\|x_0^{**}\| = 1$. Por el Teorema de Hahn-Banach, existe $x^{**} \in X^{**}$ con $\|x^{**}\| = 1$ tal que $x^{**}|_M = x_0^{**}$. En particular, se tiene que:
                \begin{equation*}
                    x^{**}(\varphi) = x_0^{**}(\varphi) = \|\varphi\|.
                \end{equation*}


                \item[Opción 2:] Consideramos $B=\{0\}\subset X^*$, que es cerrado. Como $\varphi\in X^*\setminus B$, existe $x^{**}\in \{0\}^\perp=X^{**}$ con $\|x^{**}\|=1$ tal que:
                \begin{equation*}
                    x^{**}(\varphi) = \mathrm{dist}(\varphi, B) = \|\varphi\|.
                \end{equation*}
            \end{description}

            De hecho, puesto que $\varphi$ no alcanza su norma, $|\varphi|$ no alcanza su supremo en $S_X$. Como sí existe $x^{**}$ que alcanza la norma de $\varphi$, se deduce que $x^{**}\in X^{**}\setminus J_X(X)$, de donde $J_X(X)\subsetneq X^{**}$ y, por lo tanto, $X$ no es reflexivo.
        \end{enumerate}
    \end{ejercicio}

    \begin{ejercicio}
        Rodee las afirmaciones verdaderas (que pueden ser una, dos, tres o cuatro):
        \begin{enumerate}[label=\alph*)]
            \item Sea $X$ un subespacio cerrado de $\ell_{3000}$ de dimensión infinita ¿Puede ser $X$ isomorfo a $\ell_1$?
            \begin{enumerate}[label={a\arabic*)}]
                \item Sí, porque $\ell_1 \subseteq \ell_{3000}$
                \item No, porque $\ell_1 \nsubseteq \ell_{3000}$
                \item No, porque $\ell_1$ es separable y $X$ no lo es
                \item No, porque $X^*$ es separable y $\ell_1^*$ no lo es
            \end{enumerate}

            Las dos primeras opciones no están relacionadas, por lo que no son ciertas. Veamos en primer lugar una forma de ver que no pueden ser isomorfos.
            \begin{itemize}
                \item Como $X$ es un subespacio cerrado de $\ell_{3000}$, que es reflexivo, se deduce que $X$ es reflexivo. Como $\ell_1$ no es reflexivo, se deduce que $X$ y $\ell_1$ no pueden ser isomorfos.
            \end{itemize}

            No obstante, esto no ayuda a resolver este apartado. Tenemos que:
            \begin{itemize}
                \item $\ell_1$ es separable.
                \item $\ell_{1}^*$ es isométricamente isomorfo a $\ell_\infty$, que no es separable.
                \item $X^*$ es isométricamente isomorfo a $\ell_{3000}^*/X^\perp$. Como $\ell_{3000}^*$ es separable y la aplicación proyección al cociente es continua y sobreyectiva, se deduce que $\ell_{3000}^*/X^\perp$ es separable. Por lo tanto, $X^*$ es separable.
                \item Como $\ell_{3000}$ es separable, entonces $X$ es separable.
            \end{itemize}

            Por tanto, por descarte, la única afirmación verdadera es la \textbf{a4)}. Esta además es verdadera pues, si $X$ y $\ell_1$ fuesen isomorfos, entonces $X^*$ y $\ell_1^*$ serían isomorfos, lo que es imposible pues $X^*$ es separable y $\ell_1^*$ no lo es.


            \item Sea $\psi \colon \ell_2 \to \bb{R}$ el funcional lineal dado por $\psi(x) = \sum\limits_{n=1}^\infty \frac{x(n)}{n}$ para $x \in \ell_2$:
            \begin{enumerate}[label={b\arabic*)}]
                \item $\psi$ alcanza su norma porque $\ell_2$ es reflexivo
                \item $\psi$ alcanza su norma porque $\ell_2$ es un espacio de Hilbert
                \item $\psi$ no alcanza su norma
                \item $\psi$ no alcanza su norma porque $\ell_2$ no es completo
            \end{enumerate}

            Como $\psi\in \ell_2^*$ y $\ell_2$ es reflexivo, se deduce que $\psi$ alcanza su norma (equivalentemente, $|\psi|$ alcanza su supremo en la bola unidad de $\ell_2$). Por tanto, opción \textbf{b1)} es verdadera.

            % // TODO: Terminar
            \item Sea $Y = \bb{K}^{(\bb{N})}$ (el espacio vectorial de las sucesiones casi nulas):
            \begin{enumerate}[label={c\arabic*)}]
                \item la norma de la suma y la norma del máximo son equivalentes en $Y$
                \item la norma de la suma y la norma del máximo no son equivalentes en $Y$
                \item cualesquiera dos normas definidas en $Y$ son equivalentes
                \item no se puede definir ninguna norma en $Y$
            \end{enumerate}

            Veamos que $\|\cdot\|_1$ y $\|\cdot\|_\infty$ no son equivalentes en $Y$. Para ello, fijado $n \in \bb{N}$, consideremos la sucesión:
            \begin{equation*}
                x_n = \sum_{k=1}^n e_k \in Y,
            \end{equation*}
            donde $e_k$ es la sucesión que tiene un 1 en la posición $k$ y 0 en el resto de posiciones. Claramente $x_n \in Y$ y se tiene que:
            \begin{equation*}
                \|x_n\|_1 = n \quad \text{y} \quad \|x_n\|_\infty = 1.
            \end{equation*}

            Por tanto, si $\|\cdot\|_1$ y $\|\cdot\|_\infty$ fueran equivalentes, existiría una constante $C > 0$ tal que:
            \begin{equation*}
                n = \|x_n\|_1 \leq C \|x_n\|_\infty = C \quad \forall n \in \bb{N}.
            \end{equation*}

            Esto es imposible, por lo que se concluye que $\|\cdot\|_1$ y $\|\cdot\|_\infty$ no son equivalentes en $Y$. Por lo tanto, la única afirmación verdadera es la \textbf{c2)}.

            \item En un espacio normado $X$ tomamos un conjunto convexo y cerrado $C$ y un punto $x_0 \notin C$:
            \begin{enumerate}[label={d\arabic*)}]
                \item existe $f \in X^*$ tal que $\sup \Re f(C) \leq \Re f(x_0)$
                \item existe $f \in X^*$ tal que $\sup \Re f(C) < \Re f(x_0)$
                \item no se puede asegurar que $C$ se pueda separar de $x_0$ si el interior de $C$ es vacío
                \item existe un único $f \in X^*$ tal que $\sup \Re f(C) \leq \Re f(x_0)$
            \end{enumerate}

            Suponemos que $C\neq \emptyset$, pues si $C=\emptyset$ no se puede considerar su separación de cualquier otro conjunto. En este caso, puesto que ambos conjuntos podrían tener interior vacío, no se puede aplicar de manera directa la separación en espacios normados. No obstante, puesto que $d(x_0, C) > 0$, se puede aplicar la separación fuerte en espacios normados. Como $\{x_0\},C$ son conjuntos no vacíos y convexos tal que $d(\{x_0\}, C) = d(x_0, C) > 0$, se deduce que existen $f \in X^*$ y $\alpha \in \bb{R}$ tales que:
            \begin{equation*}
                \sup \Re f(C) \leq  \alpha<\alpha + d(x_0, C) \leq \Re f(x_0).
            \end{equation*}

            Por lo tanto, las opciones \textbf{d1)} y \textbf{d2)} son verdaderas, y la opción \textbf{d3)} es falsa. La opción \textbf{d4)} es falsa porque no hay unicidad del funcional separador (ni siquiera exigiendo que tenga norma $1$). 

            Por ejemplo, en $X = \mathbb{R}^2$ con la norma euclídea, consideremos el conjunto convexo y cerrado $C = \{(0,0)\}$ y el punto $x_0 = (1,0) \notin C$. Los siguientes dos funcionales distintos $f_1, f_2 \in (\mathbb{R}^2)^*$, ambos de norma $1$, separan $C$ de $x_0$:
            \begin{align*}
                f_1(x,y) &= x \quad \Longrightarrow \quad \sup f_1(C) = 0 < 1 = f_1(x_0), \\
                f_2(x,y) &= \frac{1}{\sqrt{2}}(x+y) \quad \Longrightarrow \quad \sup f_2(C) = 0 < \frac{1}{\sqrt{2}} = f_2(x_0).
            \end{align*}
            \item Sea $Z$ un espacio normado de dimensión finita:
            \begin{enumerate}[label={e\arabic*)}]
                \item cualquier aplicación lineal que sale de $Z$ es abierta
                \item cualquier aplicación lineal que llega a $Z$ es continua
                \item cualquier aplicación lineal sobreyectiva que llega a $Z$ es abierta
                \item cualquier aplicación lineal que sale de $Z$ es continua
            \end{enumerate}

            Suponemos que todos los espacios vectoriales son normados, pues en necesario considerar una topología para poder hablar de continuidad y de aplicaciones abiertas. En el caso de la asignatura de Análisis Funcional, la topología empleada, salvo que se indique lo contrario, es la topología inducida por la norma.
            \begin{itemize}
                \item Consideramos el operador nulo:
                \Func{T}{Z}{X}{z}{0}

                Este operador es lineal y no es abierto, por lo que la afirmación \textbf{e1)} es falsa.

                \item Se ha visto en teoría que, si $X$ es un espacio normado de dimensión infinita, entonces existe un funcional lineal $f \in X^*$ que no es continuo. Por lo tanto, la afirmación \textbf{e2)} es falsa.
                
                Veamos cuál podría ser un ejemplo de un funcional lineal no continuo. Consideremos el espacio normado $c_{00}$, que sabemos es de dimensión infinita y tiene una base canónica $\{e_n\}_{n \in \bb{N}}$. Definimos el funcional lineal $f \colon c_{00} \to \bb{K}$ por:
                \begin{align*}
                    f(e_n) = n \quad \forall n \in \bb{N}.
                \end{align*}

                Este funcional es lineal, pero no es continuo pues, para cada $n \in \bb{N}$, se tiene que $\|e_n\|_\infty = 1$ y $|f(e_n)| = n$, por lo que no existe una constante $C > 0$ tal que $|f(x)| \leq C \|x\|_\infty$ para todo $x \in c_{00}$. Como $ \dim \bb{K}=1$ finita, se deduce que $f$ es un funcional lineal no continuo que llega a un espacio de dimensión finita. Por lo tanto, la afirmación \textbf{e2)} es falsa.
                
                \item Dado un operador lineal $T \colon X\to Z$ con $Z$ de dimensión finita sobreyectivo, en la Relación de Ejercicios se demostró que es abierto. Por lo tanto, la afirmación \textbf{e3)} es verdadera.

                \item Por un corolario del Teorema de Hausdorff, todo operador lineal entre espacios normados en el que el dominio es de dimensión finita es continuo. Por lo tanto, la afirmación \textbf{e4)} es verdadera.
            \end{itemize}
        \end{enumerate}
    \end{ejercicio}


\end{document}